% Long-form draft (full version, used as source for arXiv + workshop trims)
% "Prediction Bottlenecks Don't Discover Causal Structure (But Here's What
%  They Do)"

\documentclass{article}
\usepackage[preprint]{../neurips_2026}
\usepackage[utf8]{inputenc}
\usepackage[T1]{fontenc}
\usepackage{hyperref}
\usepackage{url}
\usepackage{booktabs}
\usepackage{amsfonts}
\usepackage{microtype}
\usepackage{graphicx}
\usepackage{xcolor}
\usepackage{amsmath}
\usepackage{amssymb}
\usepackage{multirow}

\title{Prediction Bottlenecks Don't Discover Causal Structure\\
(But Here's What They Actually Do)}

\author{%
  Ankit Hemant Lade\\
  \texttt{ankit@example.com}\\
  \And
  Sai Krishna\\
  \texttt{sai@example.com}\\
  \And
  Aman Chadha\\
  \texttt{aman@example.com}\\
}

\begin{document}
\maketitle

\begin{abstract}
A Mamba state-space model trained only for next-step prediction appears,
at first glance, to recover Granger-causal structure through a strikingly
simple readout: $\mathbf{S} = |\mathbf{W}_{\text{out}} \mathbf{W}_{\text{in}}|$.
Early experiments suggested this worked across architectures, converged
fast, and benefited disproportionately from interventional data with
$p < 10^{-5}$. We report the result of walking this claim through a
staged falsification pipeline covering five hypothesis layers
(architecture, sparse-regression parity, real-data recovery, intervention
amplification, and method-specificity). The method-level claim does not
survive: (i) a plain linear bottleneck does it as well or better than
the SSM, (ii) tuned Lasso beats the bottleneck on graph recovery across
synthetic and real data (NOAA climate, sector ETFs), and (iii) the
headline intervention advantage was $\sim$60\% a size-matched-data
confound plus a target-corruption artifact; when proper do(X=c)
interventions replace the random-forcing scheme, the effect vanishes,
and classical Granger shows a \emph{larger} size-matched intervention
gain than the bottleneck. We document what survives as a narrow
characterization result: bottlenecks and sparse regression behave
similarly on lagged dependency recovery, with bottlenecks marginally
preferred in mild-nonlinearity and target-corruption regimes. We hope
the falsification scaffold is as useful to the community as the
positive findings we could not salvage.
\end{abstract}

% ============================================================
\section{Introduction}
% ============================================================

Learned neural networks frequently contain more structure than their
training objectives explicitly demand. A forecasting model trained on
multivariate time series, for example, must implicitly route predictive
information from each input variable to each output variable through
its learned weights. A natural hope is that this routing structure
exposes causal or dependency relationships without any causal-specific
training objective — that is, \emph{causal discovery as a byproduct of
prediction}. If true, this would be a remarkably cheap and general
tool: every pretrained time-series forecaster would double as a
causal-discovery engine.

This paper reports the result of taking that hope seriously, across
a staged sequence of falsification experiments that tested five
distinct variants of the claim at progressively stronger controls. Our starting point was the
observation that for a Mamba-style selective state-space
model~\citep{gu2023mamba} trained for next-step MSE prediction, the
product $\mathbf{S} = |\mathbf{W}_{\text{out}}
\mathbf{W}_{\text{in}}|$ of the output and input projection matrices
appears to recover Granger-causal graph structure with striking
reliability. Initial experiments suggested the phenomenon generalized
across architectures, that it emerged from training rather than
random initialization, and, most provocatively, that interventional
data amplified the bottleneck's advantage over sparse regression
baselines with highly significant Wilcoxon $p$-values.

Each of these claims turned out to be less robust than it first
appeared. This paper documents the full falsification chain and what
ultimately survives. Our primary contribution is therefore
\emph{negative}: we show that a natural method candidate fails in
several subtly different ways, we identify the
specific controls that caused it to fail, and we report the narrow
characterization findings that remain valid after those controls are
applied. Our secondary contribution is a reusable experimental
scaffold — synthetic generators, baseline harnesses, and
intervention-control protocols — that future work on post-hoc causal
extraction from learned models can reuse to avoid the same mistakes
we made.

\paragraph{Contributions.}
\begin{enumerate}
  \item We replicate the empirical observation that prediction-trained
  bottleneck models produce a $|\mathbf{W}_{\text{out}}
  \mathbf{W}_{\text{in}}|$ score that correlates with Granger structure
  on synthetic VAR benchmarks (Section~\ref{sec:observation}).
  \item We falsify the claim that this effect is specific to
  selective state-space models. A linear bottleneck matches or beats
  the SSM at the same extraction under matched capacity and seeds
  (Section~\ref{sec:arch}).
  \item We falsify the claim that the bottleneck competes with tuned
  sparse lagged regression on graph recovery, on both synthetic
  benchmarks (CauseMe-style structural processes, $K=5$–$30$) and
  real observational data (NOAA climate, sector ETFs)
  (Section~\ref{sec:lasso}, Section~\ref{sec:realdata}).
  \item We show that an apparent interventional-data advantage,
  initially significant at $p < 10^{-5}$ across twelve
  $(K, T)$ configurations, dissolves under a size-matched control
  ($\sim$60\% of the effect) and is specific to a particular
  random-forcing intervention scheme ($\sim$the remainder), with the
  residual reproducing as a larger effect under classical Granger
  (Section~\ref{sec:intervention}).
  \item We document what does survive: a mild-nonlinearity regime
  where the bottleneck marginally outperforms Lasso, a
  sample-efficiency advantage under growing $T$, and robustness to
  per-step target corruption (Section~\ref{sec:survives}).
\end{enumerate}

% ============================================================
\section{Method: The Extraction We Are Falsifying}
\label{sec:observation}
% ============================================================

For a $K$-variable multivariate time series $\mathbf{X} \in
\mathbb{R}^{T \times K}$, we train a model with input projection
$\mathbf{W}_{\text{in}} \in \mathbb{R}^{d \times K}$ and output
projection $\mathbf{W}_{\text{out}} \in \mathbb{R}^{K \times d}$ on
next-step MSE prediction only. No causal-specific loss, no
group-lasso, no graph-structural prior.

After training, we extract:
\begin{equation}
  \mathbf{S} = \left| \mathbf{W}_{\text{out}} \mathbf{W}_{\text{in}} \right|
  \;\in\; \mathbb{R}^{K \times K},
  \qquad \mathbf{S}_{i,j} = \text{strength of }j \to i.
\end{equation}
We zero the diagonal, normalize by the max entry, and threshold to
obtain an estimated dependency graph. This is the \emph{entire}
method candidate.

For architectures with an explicit lagged dependence (e.g.\ Sai's
lagged bottleneck), we extend this to a lag-indexed tensor
$\mathbf{S}^{(\tau)} = |\mathbf{W}_{\text{out}}
\mathbf{W}^{(\tau)}_{\text{in}}|$. Evaluation is done against a
ground-truth binary adjacency tensor via flat-lag AUROC.

% ============================================================
\section{Architecture Does Not Matter (Falsification \#1)}
\label{sec:arch}
% ============================================================

Our first claim was that the SSM's selective state dynamics were
responsible for encoding clean dependency structure in its endpoint
projections. We ran an architecture sweep at matched parameter counts
and matched seeds over four datasets and five architectures
(Mamba SSM, Transformer, LSTM, MLP, Linear). The linear bottleneck
matches or beats the SSM on every dataset (Table~\ref{tab:arch}).
This result alone kills the ``SSM internals carry unique causal
signal'' story.

\begin{table}[t]
\centering
\small
\caption{Mean AUROC ($\pm$ std) over 10 seeds. Linear $\ge$ SSM
everywhere; MLP and LSTM substantially worse.}
\label{tab:arch}
\begin{tabular}{lccccc}
\toprule
Dataset & SSM (Mamba) & Transformer & LSTM & MLP & \textbf{Linear} \\
\midrule
VAR(1) chain $K{=}5$ & $1.00 {\pm} 0.00$ & $0.65 {\pm} 0.14$ & $0.38 {\pm} 0.16$ & $0.55 {\pm} 0.13$ & $\mathbf{1.00 \pm 0.00}$ \\
VAR(1) random $K{=}10$ & $0.77 {\pm} 0.08$ & $0.60 {\pm} 0.08$ & $0.49 {\pm} 0.08$ & $0.53 {\pm} 0.07$ & $\mathbf{0.99 \pm 0.02}$ \\
Regime switch $K{=}3$ & $0.98 {\pm} 0.08$ & $0.50 {\pm} 0.31$ & $0.53 {\pm} 0.33$ & $0.25 {\pm} 0.10$ & $\mathbf{1.00 \pm 0.00}$ \\
Lorenz $K{=}3$ & $0.52 {\pm} 0.24$ & $0.82 {\pm} 0.11$ & $0.74 {\pm} 0.25$ & $0.54 {\pm} 0.28$ & $\mathbf{0.94 \pm 0.09}$ \\
\bottomrule
\end{tabular}
\end{table}

\paragraph{Why this matters.} The linear bottleneck's
$\mathbf{W}_{\text{out}} \mathbf{W}_{\text{in}}$ is, by construction,
a low-rank linear estimator of the lag-1 regression matrix. The
``emergence'' we observed is just this estimator doing its job. The
SSM is a more complex model that achieves the same estimator but no
better.

% ============================================================
\section{Tuned Lasso Beats The Bottleneck (Falsification \#2)}
\label{sec:lasso}
% ============================================================

A natural follow-up: maybe the bottleneck outperforms classical
sparse regression in some specific regime (small samples, high
dimension, non-linearity). We ran the bottleneck head-to-head against
validation-tuned Lasso and reduced-rank regression across a focused
grid ($K \in \{10, 20\}, T \in \{150, 300\}$, max-lag 8) and then on
domain-realistic tigramite structural processes at $K \in \{10, 20,
30\}$. With untuned OLS/Ridge/Lasso baselines, the bottleneck looked
promising. With \emph{tuned} Lasso and RRR, the advantage evaporates
(Table~\ref{tab:lasso}).

\begin{table}[t]
\centering
\small
\caption{Focused stress grid (Exp 15): once baselines are tuned,
bottleneck wins only 12\% of graph-recovery runs and 0\% of
prediction-MSE runs.}
\label{tab:lasso}
\begin{tabular}{lrr}
\toprule
Metric vs.\ best tuned baseline & Mean flat-lag AUROC delta & Win rate \\
\midrule
Bottleneck vs.\ \{OLS, Ridge, Lasso, RRR\} & $-0.050$ & $0.12$ \\
Bottleneck prediction MSE & $+0.286$ & $0.00$ \\
\midrule
Best baseline by AUROC & Lasso ($38/48$), RRR ($10/48$) & \\
Best baseline by MSE & Lasso ($48/48$) & \\
\bottomrule
\end{tabular}
\end{table}

On tigramite CauseMe-style benchmarks, the gap widens with $K$:
tuned Lasso reaches AUROC 0.98 at $K=20$ while the bottleneck sits at
0.73. No dimension or sample-size regime we found flips the result.

% ============================================================
\section{Real-World Data and Benchmarks: Bottleneck Is Not Competitive}
\label{sec:realdata}
% ============================================================

We evaluate the same six methods on three datasets of progressively
cleaner ground truth (Table~\ref{tab:real}):
(i) 63 years of NOAA monthly climate indices
(ENSO, NAO, PDO, AMO, SOI, PNA) with teleconnection edges from the
climate dynamics literature~\citep{trenberth1994decadal,delworth2000observed,newman2016pdo};
(ii) 7.5 years of daily Yahoo Finance SPDR sector ETF log returns
($K{=}10$) with sector lead-lag edges from finance commentary (treated
as soft ground truth);
(iii) Lorenz-96 ($K{=}10$, $F{=}10$, chaotic) with mathematically
unambiguous ground truth (each variable is driven by its two left
and one right neighbor)~\citep{tank2021neural,nauta2019causal}.

\paragraph{Ground-truth choices matter.} On the climate data we
deliberately \emph{exclude} the ENSO$\leftrightarrow$SOI pair from
ground truth even though it is often counted in the literature. The
Southern Oscillation Index is a surface-pressure analog of the Niño
3.4 SST anomaly: they are the same signal with opposite sign, not a
causal edge. Including them inflates every method's AUROC and, as we
show below, changes the ranking: under the inflated ground truth,
bivariate Granger ties for the top spot on climate; under the
honest ground truth, \emph{Granger drops to the bottom} and the
bottleneck drops to the middle.

\paragraph{Results.} On climate the SSM bottleneck is, surprisingly,
the top method ($0.820 \pm 0.021$ over five seeds) — but this is against soft ground truth
with only 3 edges over $K{=}6$, where small changes move AUROC
substantially; we do not read this as a causal-discovery claim. On
finance, Ridge wins decisively (0.696) and all other methods are at
or below chance (including the SSM bottleneck at $0.447 \pm 0.033$), consistent with
sector-ETF returns admitting no straightforward lagged causal
structure. On Lorenz-96, where the ground truth is unambiguous,
\emph{PCMCI and bivariate Granger dominate} (0.986, 0.979), tuned
Lasso and Ridge essentially match them (0.974), and the lagged
bottleneck trails at 0.916 while the SSM bottleneck is worst at
$0.722 \pm 0.031$. The bottleneck is not the best method on any dataset with
clean ground truth.

\begin{table}[t]
\centering
\small
\caption{AUROC on three real/benchmark datasets. The Weight-Proj (SSM)
row reports mean $\pm$ std over 5 training seeds; the lagged bottleneck
row uses a single seed. The bottleneck tops only climate, where ground
truth is weak (3 edges, $K{=}6$) and AUROC is sensitive to individual
edge choices; on the clean Lorenz-96 benchmark, classical causal
methods lead by a wide margin.}
\label{tab:real}
\begin{tabular}{lccc}
\toprule
Method & Climate (NOAA) & Finance (ETFs) & Lorenz-96 \\
 & $K{=}6, T{=}757$ & $K{=}10, T{=}1893$ & $K{=}10, T{=}1500$ \\
\midrule
Weight-Proj (SSM, 5 seeds) & $\mathbf{0.820 \pm 0.021}$ & $0.447 \pm 0.033$ & $0.722 \pm 0.031$ \\
PCMCI & $0.704$ & $0.510$ & $\mathbf{0.986}$ \\
Tuned Lasso & $0.630$ & $0.500$ & $0.974$ \\
Lagged Bottleneck & $0.617$ & $0.484$ & $0.916$ \\
Tuned Ridge & $0.617$ & $\mathbf{0.696}$ & $0.974$ \\
Bivariate Granger & $0.605$ & $0.498$ & $0.979$ \\
\bottomrule
\end{tabular}
\end{table}

% ============================================================
\section{Interventions: A Confound Story}
\label{sec:intervention}
% ============================================================

Advisor feedback on an earlier version of this work suggested that if
observational data yielded parity, \emph{interventional} data might
amplify differences between learned bottlenecks and sparse
regression. We took this seriously. The initial experiment (Exp 19)
compared $\text{AUROC}(\mathbf{X}_{\text{obs}})$ against
$\text{AUROC}(\mathbf{X}_{\text{obs}} \cup \mathbf{X}_{\text{int}})$
for both bottleneck and Lasso, on four graph types.

We reported $12 / 12$ $(K, T)$ configs significant at
$p < 10^{-5}$ in favor of the bottleneck, with win rates of
$80\%$–$100\%$. This looked like our one real positive result.

It does not survive two controls.

\paragraph{Control 1 (Exp 22): sample size.}
Exp 19 compared an observational set of size $T$ against a combined
set of size $T + K \cdot T_{\text{int/target}}$. The
``intervention gain'' conflated \emph{more data} with
\emph{intervention content}. Exp 22 adds a third arm,
$\mathbf{X}_{\text{obs,big}}$, a fresh observational draw of the same
total size as the combined set. The intervention-specific gain is
then $\text{AUROC}(\text{combined}) -
\text{AUROC}(\text{obs,big})$. The resulting size-matched gain is
about $\sim$40\% of the raw effect.

\paragraph{Control 2 (Exp 23): intervention type.}
Exp 21's ``intervention'' was a per-step random forcing
$x_{i,t} \leftarrow s \cdot \epsilon_{i,t}$, which is not a standard
$\text{do}$-intervention. Exp 23 swaps in proper constant clamps
$\text{do}(X_i = c)$ and soft (added-noise) interventions, and adds
Ridge and RRR baselines. Results are in Table~\ref{tab:int_type}.
At $K=30$ with constant interventions, Ridge \emph{beats} the
bottleneck (size-matched gap $-0.021$, $p = 1.00$). Only the
random-forcing scheme produces a stable bottleneck gain, and that
gain is best explained as robustness to corrupted target rows, not
as causal-structure exposure.

\begin{table}[t]
\centering
\small
\caption{Intervention-specific gap vs.\ the best tuned baseline,
by intervention type (mean over $K \in \{20, 30\}$, 12 seeds).}
\label{tab:int_type}
\begin{tabular}{lcc}
\toprule
Intervention type & Significant configs & Mean gap (AUROC) \\
\midrule
\texttt{random} (Exp 21 scheme) & $12/12$ & $+0.116$ \\
\texttt{constant} ($\text{do}(X_i = c)$) & $3/12$ & $+0.002$ \\
\texttt{soft} (exogenous noise) & $6/12$ & $+0.008$ \\
\bottomrule
\end{tabular}
\end{table}

\paragraph{Control 3 (Exp 25): method-agnosticity.}
The residual effect is not unique to the bottleneck. When we add
PCMCI and bivariate Granger to the same size-matched control, the
bottleneck shows a significant size-matched intervention gain
($+0.026$ to $+0.054$, $p < 10^{-3}$ across $K \in \{10, 20, 30\}$),
\emph{but so does Granger}, with larger effect sizes at higher $K$
($+0.040$ at $K=20$, $+0.095$ at $K=30$, same significance).
Lasso uniquely does not benefit and is hurt at low $K$, consistent
with per-equation regression breaking under target corruption.
PCMCI hovers near random. Bottom line: the intervention effect,
insofar as it survives controls, is method-agnostic and matches a
target-corruption-robustness story rather than a
causal-structure-extraction one.

% ============================================================
\section{What Survives}
\label{sec:survives}
% ============================================================

After five falsifications, three narrow positives remain:

\begin{enumerate}
  \item \textbf{Mild-nonlinearity regime.} In a regime sweep over
  noise, density, and nonlinearity strength (Exp 18), the bottleneck
  outperforms Lasso at \texttt{nonlinear=0.3} (78\% win rate over 3
  seeds and 9 hyperparameter cells). At stronger nonlinearity,
  both methods fail equally.
  \item \textbf{Sample-efficiency under growing $T$.} Bottleneck
  gains from additional observational data exceed Lasso's by
  $\sim 0.07$ AUROC at $K \in \{20, 30\}$ (Exp 22, size-matched
  control arms). This is a useful but modest result for settings
  where collecting more observations is cheap.
  \item \textbf{Target-corruption robustness.} The bottleneck's
  shared $\mathbf{W}_{\text{out}}$ absorbs per-step random
  forcing better than Lasso's per-equation fits (Exp 21--23, 25).
  This explains the residual intervention effect but is not a causal
  discovery result — it is a reliability result.
\end{enumerate}

These are characterization claims, not method claims. We do not
propose the bottleneck as a causal discovery method.

% ============================================================
\section{Lessons}
% ============================================================

Three methodological points stand out from this falsification chain,
beyond the specific negative results themselves.

\paragraph{Size-matched controls matter more than seed counts.}
Exp 19 had 5 seeds. Exp 21 had 20 seeds and reported $p < 10^{-5}$
across 12 configs. Both found the same effect. Neither would have
told us that the effect was largely a sample-size confound. Adding
seeds tightens a noisy measurement; it cannot diagnose a structural
control that the comparison omits.

\paragraph{Pre-register the intervention scheme.}
The gap between ``random forcing'' (per-step clamp to
$s \cdot \epsilon$) and proper $\text{do}(X_i = c)$ is the
difference between a significant effect and a null, and the gap is
not obvious from the code. We should have written down which
intervention semantics counted as ``interventional data'' before
running the experiment.

\paragraph{Prediction baselines are not causal baselines.}
Exps 15–23 compared the bottleneck only against prediction-fit
baselines (OLS, Ridge, Lasso, RRR). It was natural to miss that
classical \emph{causal} discovery methods (Granger, PCMCI) would
benefit from the same intervention scheme by the same mechanism.
Exp 25 added them and killed the method-specificity story. Any
future ``neural causal discovery'' claim should include at least
one classical causal baseline in the intervention control arm.

% ============================================================
\section{Conclusion}
% ============================================================

We set out to find a free causal discovery method hiding inside
prediction bottlenecks and could not. The headline effects that
initially looked promising either do not replicate under size-matched
controls, are specific to a non-standard intervention scheme, or are
method-agnostic once real causal baselines are tested. What survives
is a narrow set of characterization findings about when bottleneck
and sparse regression behave similarly versus when they diverge, and
a reusable experimental scaffold for others who want to try the same
project without repeating the same confounds.

We have deliberately left the positive framing out of the title.
The contribution is the falsification, honestly reported.

\paragraph{Code.} All experiments, data loaders, and a one-line
reproduction script are available at
\texttt{https://github.com/ankitlade12/ssm-causal}.

\bibliography{../references}

\end{document}
